\section*{Chapitre \ref{ch:dv:no-arbitrage-and-the-fundamental-theorems} --- Absence d'arbitrage et théorèmes fondamentaux}

\begin{solution}{exo:dv:no-arbitrage-and-the-fundamental-theorems:1}
$q_u+q_d=0.99$ et $115q_u+95q_d=100$ donnent $20q_u=100-94.05$~: $q_u=0.2975$,
$q_d=0.6925$. La probabilité risque-neutre de hausse est $0.2975/0.99=0.3005$~; le call
paie $(15,0)$ et vaut $15q_u=4.4625$.
\end{solution}

\begin{solution}{exo:dv:no-arbitrage-and-the-fundamental-theorems:2}
En dessous de 5\,000~: les calls ne paient rien, le put long à 6\,000 paie $6000-S$, le
put court à 5\,000 coûte $5000-S$~; total 1\,000. Entre les deux~: le call long paie $S-5000$, le
put long $6000-S$~; total 1\,000. Au-dessus de 6\,000~: les calls paient $(S-5000)-(S-6000)=1000$,
les puts ne paient rien.
\end{solution}

\begin{solution}{exo:dv:no-arbitrage-and-the-fundamental-theorems:3}
Le put paie $(0,0,20)$~; son intervalle est $(20\times0,\ 20\times0.44)=(0,\ 8.80)$,
donc 9.00 constitue un \omterm{def:dv:no-arbitrage-and-the-fundamental-theorems:arbitrage}{arbitrage}. Vendre le put à 9.00 et le sur-répliquer avec le
portefeuille qui l'égale dans les états de hausse et de baisse~: $-0.5$ action et 60 obligations,
qui paie $(0,10,20)$ et coûte $-50+58.8=8.80$. L'edge vaut 0.20 aujourd'hui, plus 10 dans
l'état intermédiaire.
\end{solution}

\begin{solution}{exo:dv:no-arbitrage-and-the-fundamental-theorems:4}
Une obligation qui paie $(1,1)$ à 0.95 et un actif qui paie $(1,0)$ au prix de 0. La matrice des
paiements est inversible, donc chaque paiement possède un unique \omterm{def:dv:no-arbitrage-and-the-fundamental-theorems:claim}{portefeuille de réplication} et un seul coût~:
la \omterm{def:dv:no-arbitrage-and-the-fundamental-theorems:loop}{loi du prix unique} est vérifiée. Mais acheter le second actif ne coûte rien et paie 1
dans le premier état~: un \omterm{def:dv:no-arbitrage-and-the-fundamental-theorems:arbitrage}{arbitrage}. Le vecteur unique de \omterm{def:dv:no-arbitrage-and-the-fundamental-theorems:ad}{prix d'état} est $(0,0.95)$, qui
n'est pas strictement positif.
\end{solution}

\begin{solution}{exo:dv:no-arbitrage-and-the-fundamental-theorems:5}
Les probabilités risque-neutres sont $q/0.98$, soit 0.3061, 0.4898 et 0.2041~;
la \omterm{def:dv:no-arbitrage-and-the-fundamental-theorems:spd}{densité des prix d'état} est $q/\P$, soit 0.6667, 1.0667 et 2.0000. Elle est la plus élevée dans l'état de baisse~: une unité versée lorsque
l'action a chuté vaut deux fois sa valeur pondérée par les probabilités, soit le prix
de l'assurance contre les mauvais états.
\end{solution}

\begin{solution}{exo:dv:no-arbitrage-and-the-fundamental-theorems:6}
Elle vaut $q_u$, strictement compris entre 0.10 et 0.54. À 0.54, elle est sur-répliquée
par le portefeuille qui paie $(1,0,0)$ dans les états de hausse et de baisse~: $\theta_s=1/40=0.025$
action et $-2$ obligations, qui paie $(1,0.5,0)$ et coûte $2.5-1.96=0.54$.
\end{solution}

\begin{solution}{exo:dv:no-arbitrage-and-the-fundamental-theorems:7}
Les deux sommets donnent 10.80~: l'intervalle se réduit au seul point 10.80. Le spread paie
$(20,10,0)$, soit $0.5$ action moins 40 obligations dans chaque état~: il est réplicable,
son prix est donc imposé bien que le marché soit incomplet.
\end{solution}

\begin{solution}{exo:dv:no-arbitrage-and-the-fundamental-theorems:8}
À l'intérieur de l'intervalle, aucun prix n'est un \omterm{def:dv:no-arbitrage-and-the-fundamental-theorems:arbitrage}{arbitrage} et aucun n'est «~équitable~» sans un
choix de \omterm{def:dv:no-arbitrage-and-the-fundamental-theorems:ad}{prix d'état}~: 6 est le choix du modèle, non celui du marché. Acheté à 5, le
call ne peut pas être couvert pour garantir un profit certain, puisque le marché ne peut pas le répliquer~; la
meilleure couverture statique sous-réplique le call pour 2 et laisse le risque de l'état-$u$. Le profit
de 1 n'existe que si les \omterm{def:dv:no-arbitrage-and-the-fundamental-theorems:ad}{prix d'état} du modèle sont justes, ce qui relève du pari.
\end{solution}

\begin{solution}{pb:dv:no-arbitrage-and-the-fundamental-theorems:1}
\textbf{1.} Call long à 5\,000 $(S-5000)^+$, call court à 6\,000 $-(S-6000)^+$, put court
à 5\,000 $-(5000-S)^+$, put long à 6\,000 $(6000-S)^+$~; la somme des calls donne
$\min(\max(S-5000,0),1000)$ et celle des puts $1000-$ ce montant~: 1\,000.
\textbf{2.} 100\,000~USD.
\textbf{3.} $973.40-324.60-156.70+470.00=962.10$ points, soit 96\,210~USD.
\textbf{4.} $971.80-326.20-158.30+468.40=955.70$ points, soit 95\,570~USD.
\textbf{5.} Elle élimine tout le risque d'indice~; ce qui subsiste est la contrepartie (la
chambre de compensation), le risque opérationnel d'une jambe exécutée sans les autres, et
le financement de la marge sur un box vendeur.
\textbf{6.} $-\ln(0.96210)=3.86\%$.
\textbf{7.} $-\ln(0.95570)=4.53\%$.
\textbf{8.} $958.90$~; $-\ln(0.95890)=4.20\%$.
\textbf{9.} $1000-962.10=37.90$ points, soit 3\,790~USD.
\textbf{10.} Au prêt~: $-13.6$~pb, à l'emprunt~: $+53.1$~pb.
\textbf{11.} Exécuter les jambes une à une coûte la moitié de la largeur sur chacune des quatre cotations (6.40 points en
tout ici)~; un package se traite à un seul prix proche du prix milieu et supprime le
risque d'exécution partielle.
\textbf{12.} Le box juste vaut $1000e^{-0.04}=960.79$. L'achat coûte 962.10, plus que
cela, et la vente rapporte 955.70, moins~: pas d'\omterm{def:dv:no-arbitrage-and-the-fundamental-theorems:arbitrage}{arbitrage} (\texttt{check\_box} renvoie
rien).
\textbf{13.} 960.79 points.
\textbf{14.} Les jambes courtes pourraient être assignées de façon anticipée, transformant le box en
position ouverte~; le paiement ne serait plus certain.
\textbf{15.} Les bons du Trésor offrent un rendement d'opportunité (convenience yield), tiré de leur usage
comme collatéral et de leur liquidité, d'environ 40~pb en moyenne dans l'étude citée, si bien que leur rendement est inférieur au
taux sans risque que révèlent les box.
\textbf{16.} Les détenteurs à effet de levier de portefeuilles indiciels, les fonds et les teneurs de marché qui
souhaitent emprunter sur marge à un taux proche du taux sans risque plutôt
que le taux de marge d'un courtier.
\textbf{17.} $10\,000\,000/95\,890=104.3$~: 104 box, versant 9\,972\,560~USD aujourd'hui et
recevant 10\,400\,000~USD dans un an.
\textbf{18.} Les positions sont compensées à la chambre de compensation, qui les transfère à un autre
membre compensateur ou les clôture~; sa créance porte sur la chambre de compensation, dont les
ressources en cas de défaut la garantissent.
\textbf{19.} 4.20~\%, 19.7~pb au-dessus du taux overnight de 4.00~\%.
\textbf{20.} Parce qu'il verse l'écart entre les strikes dans chaque état, et un montant certain
à $T$ est exactement ce que verse une obligation zéro-coupon.
\end{solution}

\begin{solution}{iq:dv:no-arbitrage-and-the-fundamental-theorems:1}
Call long et put court en $K_1$, call court et put long en $K_2$~: deux
forwards synthétiques qui se compensent en le montant certain $K_2-K_1$ à l'échéance. Son prix vaut
$P(0,T)(K_2-K_1)$, il cote donc un taux sans risque pour l'échéance de l'option, avec la
chambre de compensation comme contrepartie~; il nécessite un exercice européen.
\iqlookfor{parité call-put appliquée deux fois, et taux lu sur le prix.}
\end{solution}

\begin{solution}{iq:dv:no-arbitrage-and-the-fundamental-theorems:2}
L'absence d'\omterm{def:dv:no-arbitrage-and-the-fundamental-theorems:arbitrage}{arbitrage} équivaut à l'existence d'une probabilité équivalente à la probabilité réelle
pour laquelle les prix actualisés des actifs échangés sont des martingales (en temps continu~:
absence de free lunch avec risque évanescent, et une sigma-martingale). Ses probabilités sont
des \omterm{def:dv:no-arbitrage-and-the-fundamental-theorems:ad}{prix d'état} normalisés~: elles codent des prix, primes de risque incluses, et non des croyances
sur la vraisemblance des états.
\iqlookfor{l'équivalence, et «~pondérations de prix, non prévisions~».}
\end{solution}

\begin{solution}{iq:dv:no-arbitrage-and-the-fundamental-theorems:3}
Avec un taux nul, $110q_u+95(1-q_u)=100$ donne $q_u=1/3$~: le digital vaut 0.3333.
\iqlookfor{résoudre pour la probabilité risque-neutre, et non utiliser une probabilité réelle.}
\end{solution}

\begin{solution}{iq:dv:no-arbitrage-and-the-fundamental-theorems:4}
Les vecteurs de \omterm{def:dv:no-arbitrage-and-the-fundamental-theorems:ad}{prix d'état} résolvent $D^\top q=p$~; la solution est unique exactement lorsque les
gains engendrent chaque état, ce qui définit la complétude. Un marché trinomial avec une action
et une obligation est incomplet~: les prix d'options ne sont que bornés, et tout vecteur positif
de \omterm{def:dv:no-arbitrage-and-the-fundamental-theorems:ad}{prix d'état} fournit un prix sans \omterm{def:dv:no-arbitrage-and-the-fundamental-theorems:arbitrage}{arbitrage}~; en choisir un relève d'un modèle.
\iqlookfor{algèbre linéaire sous-jacente au théorème, et bornes au lieu de prix.}
\end{solution}

\begin{solution}{iq:dv:no-arbitrage-and-the-fundamental-theorems:5}
Aux prix milieu, le butterfly coûte $12.00-13.00+0.90=-0.10$, soit une violation de convexité. Aux
prix exécutables, il coûte $12.05-2\times6.45+0.95=0.10$~: pas d'\omterm{def:dv:no-arbitrage-and-the-fundamental-theorems:arbitrage}{arbitrage}. Les
prix milieu ne sont pas des prix négociables~; la question est de savoir si les cotations le sont.
\iqlookfor{vérifier au cours acheteur et au cours vendeur, et non au prix milieu.}
\end{solution}

\begin{solution}{iq:dv:no-arbitrage-and-the-fundamental-theorems:6}
Gardez la chaîne triée par échéance et strike, avec le forward et le facteur d'actualisation par
échéance~; à chaque mise à jour de cotation, ne revérifier que les inégalités qui font intervenir le
strike modifié (bornes, voisins pour la monotonie et la pente, triplets pour la
convexité, les \omterm{def:dv:no-arbitrage-and-the-fundamental-theorems:box}{box spreads} qui le contiennent)~: un nombre constant de vérifications. Utiliser des
prix exécutables, des tolérances d'un tick, et filtrer les cotations périmées~; journaliser les violations avec
le portefeuille~; lancer périodiquement une vérification complète par programme linéaire.
\iqlookfor{vérifications incrémentales locales en strike, prix exécutables, et un
passage global périodique.}
\end{solution}
